% !TeX encoding = UTF-8
% Fuente editorial del PDF y del cuadernillo web. HVT, 10 de octubre de 2026.
\documentclass{../plantilla/hvtbook}
\usepackage{amsmath}
\pagestyle{plain}
\hypersetup{pdftitle={Hábitat regenerativo},pdfauthor={Hidrógeno Verde Turquesa}}
\HVTTitulo{Hábitat regenerativo}
\HVTSubtitulo{Una casa unifamiliar de inspiración colonial, integrada al suelo, al agua y a la energía de su entorno. Patio, corredores, vegetación y confort para la vida familiar.}
\HVTNivel{Arquitectura e ingeniería del hábitat}
\HVTSerie{Publicación técnica}
\HVTNumero{VC-01}
\HVTAccion{Explorar capítulos}
\HVTDescripcion{Hábitat regenerativo: casa campestre de inspiración colonial con patio, suelo vivo, confort, reúso de agua y autonomía energética en evaluación.}
\HVTCanonical{https://hidrogenoverdeturquesa.com/proyecto-casa-campestre}
\HVTRegreso{Volver a proyectos}{/\#portfolio}
\HVTPortada{../../images/portfolio/casa-campestre.jpg}{Muros claros, teja de barro y patio. Petra Nesti / Pexels. Referencia, no obra de HVT.}
\begin{document}
\HVTMakeTitle
\begin{HVTPage}{Contenido}{VC-01 · Versión 1}{Una casa pensada desde su lugar}
\HVTLead{Diseñar una vivienda campestre para una familia, con menor demanda de recursos, espacios confortables y autonomía evaluada mediante datos.}
\HVTText{Proyecto en formulación. Todavía deben definirse el predio, la familia usuaria, el área construida y el presupuesto. La sostenibilidad integral es el propósito del diseño; su desempeño se comprobará mediante estudios, simulaciones y seguimiento de la vivienda. Esta edición no presenta una casa construida ni ahorros medidos.}
\begin{itemize}
\item 1. Alcance y necesidades de la familia.
\item 2. Terreno, suelo y agua subterránea.
\item 3. Humedad del suelo y paisaje productivo.
\item 4. Orientación y diseño bioclimático.
\item 5. Confort y calidad del ambiente interior.
\item 6. Agua lluvia y abastecimiento.
\item 7. Tratamiento y reúso de aguas grises.
\item 8. Autonomía energética y almacenamiento.
\item 9. Alimentos, orgánicos y materiales.
\item 10. Instrumentación y evaluación del diseño.
\item 11. Etapas, responsables y entregables.
\item 12. Viabilidad, decisiones y créditos.
\end{itemize}
\end{HVTPage}
\begin{HVTPage}{Lámina}{Identidad arquitectónica}{Una casa colonial que mira hacia el patio}
\HVTImagen{../../images/portfolio/patio-colonial.jpg}{Patio de una hacienda mexicana. Amar Preciado / Pexels. Referencia espacial, no diseño ni obra de HVT.}
\HVTText{La propuesta toma como lenguaje muros de acabado claro, cubierta de teja de barro, aleros, carpintería de madera, corredores y un patio con vegetación. Los arcos y balcones se incorporarán cuando aporten a la composición y al uso, sin imponerlos a todos los espacios.}
\HVTText{El patio organizará encuentros, luz y relación con el jardín. La profundidad de los corredores, las ventanas y la cubierta se evaluarán con el clima del predio. La imagen colonial convivirá con instalaciones accesibles para mantener el agua, la energía y el monitoreo; su desempeño se comprobará, no se dará por supuesto por su apariencia.}
\end{HVTPage}
\begin{HVTPage}{Capítulo 1}{Programa familiar}{El hogar como unidad de diseño}
\HVTLead{La arquitectura parte de la vida cotidiana: habitar, descansar, trabajar, cultivar y desplazarse. Cada sistema debe ser comprensible y mantenible por quienes viven allí.}
\HVTText{La distribución estudiará habitaciones vinculadas al patio y corredores hacia el jardín. El programa se construirá con la familia: número de habitantes, horarios, accesibilidad, trabajo en casa, necesidades térmicas, cocina, lavado, almacenamiento y movilidad. Se separarán las necesidades esenciales de las ampliaciones futuras. La infraestructura técnica tendrá acceso de mantenimiento sin invadir las áreas habitables.}
\begin{HVTTable}{Decisión}{Información necesaria}{Resultado esperado}
\HVTTableRow{Vivienda}{Ocupación, actividades, privacidad y crecimiento familiar.}{Programa de espacios y relaciones, sin fijar todavía un área.}
\HVTTableRow{Predio}{Ubicación, pendientes, accesos, servicios y restricciones.}{Área apta para implantar y zonas que deben conservarse.}
\HVTTableRow{Autonomía}{Cargas esenciales y duración de interrupciones a cubrir.}{Nivel de servicio que podrá verificarse.}
\HVTTableRow{Operación}{Disponibilidad de la familia y apoyo técnico local.}{Plan de mantenimiento y costos de reposición.}
\end{HVTTable}
\HVTText{Se propone comparar una solución básica eficiente con alternativas de mayor autonomía. La vivienda, el huerto y la movilidad compartirán un balance de recursos, pero sus consumos se medirán por separado. El objetivo es tomar decisiones útiles para la familia, no sumar tecnologías sin una necesidad definida.}
\end{HVTPage}
\begin{HVTPage}{Capítulo 2}{Lectura del terreno}{Antes de dibujar la planta, conocer el suelo}
\HVTLead{El suelo cumple dos funciones distintas: soporta la construcción y sostiene la vegetación. Cada una requiere su propia caracterización.}
\HVTText{La propuesta de campaña incluye levantamiento topográfico, reconocimiento de escorrentías, observación de erosión y registro de zonas encharcables. El estudio geotécnico deberá definir las condiciones de cimentación y estabilidad; la evaluación agrícola caracterizará textura, profundidad efectiva, infiltración y materia orgánica. Un sensor de humedad no sustituye el estudio geotécnico.}
\begin{HVTTable}{Escala}{Registro propuesto}{Decisión de diseño}
\HVTTableRow{Predio}{Pendientes, cobertura, recorridos del agua y accesos.}{Implantación, conservación y drenaje superficial.}
\HVTTableRow{Construcción}{Estratigrafía, resistencia y presencia de agua.}{Cimentación y manejo de humedad bajo asesoría geotécnica.}
\HVTTableRow{Huerto}{Textura, raíces, infiltración y retención de agua.}{Especies, cobertura y sectorización del riego.}
\end{HVTTable}
\HVTText{Se evitará decidir un sistema de infiltración de aguas tratadas antes de conocer la capacidad del terreno y los receptores cercanos. Los hallazgos se incorporarán en un plano único de restricciones y oportunidades. La profundidad de exploración y los ensayos no se prescriben aquí: dependen del lugar y del profesional responsable.}
\end{HVTPage}
\begin{HVTPage}{Capítulo 3}{Suelo vivo}{Medir humedad para manejar el agua}
\HVTLead{El riego se ajustará al suelo, al cultivo y al clima, en lugar de depender únicamente de un horario fijo.}
\HVTText{La FAO distingue infiltración, capacidad de campo y agua disponible para las plantas. Estas propiedades varían con el perfil del suelo; por eso una lectura aislada de humedad no determina por sí sola cuánto regar. \HVTCite{2}}
\HVTText{HVT propone instalar puntos de observación en sectores representativos del huerto, acompañados de registros de lluvia y riego. La ubicación y profundidad se acordarán según las raíces de las especies elegidas. Los sensores se contrastarán con muestras locales y se documentarán su deriva, incertidumbre y respuesta a la temperatura.}
\begin{HVTTable}{Variable}{Cómo se registrará}{Uso en el proyecto}
\HVTTableRow{Humedad}{Sondas contrastadas con muestras del suelo.}{Detectar tendencias de secado y exceso de agua.}
\HVTTableRow{Entrada}{Lluvia y volumen aplicado por sector.}{Relacionar el riego con la respuesta observada.}
\HVTTableRow{Vegetación}{Estado del cultivo y cobertura del suelo.}{Verificar que el ahorro de agua no perjudique las plantas.}
\end{HVTTable}
\HVTText{El piloto comparará sectores equivalentes con y sin cobertura orgánica, manteniendo registro de diferencias iniciales. Se evaluarán agua aplicada y condición del cultivo durante un periodo acordado. Los umbrales automáticos se definirán después de la calibración, con operación manual disponible ante fallos.}
\end{HVTPage}
\begin{HVTPage}{Capítulo 4}{Sol y viento}{Orientar la casa según su clima}
\HVTLead{La orientación se decidirá con el recorrido solar, las sombras, el viento y los horarios de ocupación del lugar.}
\HVTText{El diseño solar pasivo relaciona implantación, ventanas, protección solar, masa térmica y distribución interior. La guía del Departamento de Energía sirve como fundamento; sus recomendaciones para otras latitudes no se trasladan automáticamente a un predio colombiano. \HVTCite{1}}
\HVTText{Se propone estudiar al menos tres alternativas de implantación sobre la misma topografía. Cada alternativa mantendrá un programa familiar equivalente para comparar soleamiento, iluminación natural, exposición al viento y demanda térmica. La selección considerará también vistas, privacidad, accesibilidad y movimiento de tierras.}
\begin{HVTTable}{Elemento}{Alternativas a comparar}{Evidencia para decidir}
\HVTTableRow{Volumen}{Orientación y distribución de espacios.}{Simulación solar y demanda por zona.}
\HVTTableRow{Envolvente}{Cubierta, aislamiento, ventanas y aleros.}{Confort horario, iluminación y costo.}
\HVTTableRow{Ventilación}{Aberturas y recorridos del aire.}{Renovación y confort según clima exterior.}
\end{HVTTable}
\HVTText{La cubierta de teja y sus aleros se coordinarán con captación de lluvia y paneles solares. Se estudiarán faldones o estructuras complementarias que conserven la lectura arquitectónica y el acceso de mantenimiento. No se fijarán una orientación universal ni un porcentaje de ventana sin evaluar el clima y el uso de cada habitación.}
\end{HVTPage}
\begin{HVTPage}{Capítulo 5}{Habitar con confort}{Comprobar cómo se siente y funciona la casa}
\HVTLead{Una vivienda eficiente debe ofrecer condiciones adecuadas para descansar y realizar las actividades cotidianas.}
\HVTText{La campaña propuesta registrará temperatura y humedad interiores, condiciones exteriores, ocupación y funcionamiento de ventanas o equipos. Se incorporarán encuestas breves de percepción de la familia y una evaluación profesional del confort, escogiendo el método según la ventilación y el clima del edificio.}
\begin{HVTTable}{Aspecto}{Comprobación propuesta}{Criterio por acordar}
\HVTTableRow{Térmico}{Medición por espacios y simulación horaria.}{Horas ocupadas dentro del rango de confort adoptado.}
\HVTTableRow{Aire}{Ventilación, fuentes interiores y seguimiento de CO2.}{Renovación adecuada; el CO2 no sustituye medir otros contaminantes.}
\HVTTableRow{Luz}{Iluminancia y observación de deslumbramiento.}{Luz útil para cada actividad sin molestias.}
\HVTTableRow{Humedad}{Inspección de filtraciones y puntos fríos.}{Ausencia de humedad persistente y causas identificadas.}
\end{HVTTable}
\HVTText{La comparación entre alternativas registrará simultáneamente confort y consumo: ahorrar energía reduciendo el servicio a la familia no será un resultado satisfactorio. Las metas se fijarán antes de seleccionar la solución definitiva y se revisarán con datos de temporadas contrastantes.}
\end{HVTPage}
\begin{HVTPage}{Capítulo 6}{Agua lluvia}{Diseñar con un balance, no solo con un tanque}
\HVTLead{La captación se dimensionará con la lluvia disponible, la superficie efectiva de cubierta y las necesidades reales de la vivienda.}
\HVTEquation{V = P A C}{Estimación preliminar: V en litros, lluvia P en milímetros, área A en metros cuadrados y coeficiente efectivo C sin unidades.}
\HVTText{Ejemplo didáctico: 50 mm de lluvia sobre 100 metros cuadrados, con un coeficiente efectivo de 0,8, representarían 4.000 litros captables antes de limitar por almacenamiento. No es una previsión para el predio. El cálculo definitivo debe representar la secuencia temporal de lluvia, descartes, pérdidas, demanda y reboses.}
\HVTText{Se propone un balance diario de volumen almacenado y una prueba de sensibilidad frente a periodos secos. Las primeras escorrentías, el material de cubierta, la limpieza, la protección del depósito y el tratamiento se definirán según calidad y uso. No se supondrá que el agua lluvia es potable.}
\begin{HVTTable}{Pregunta}{Dato de entrada}{Salida de diseño}
\HVTTableRow{Disponibilidad}{Serie de lluvia y área efectiva.}{Volumen aprovechable y variación temporal.}
\HVTTableRow{Servicio}{Demanda por uso y reserva acordada.}{Capacidad de almacenamiento y respaldo.}
\HVTTableRow{Calidad}{Caracterización y uso previsto.}{Tratamiento, controles y mantenimiento.}
\end{HVTTable}
\end{HVTPage}
\begin{HVTPage}{Capítulo 7}{Aguas grises}{Separar, tratar y verificar antes del reúso}
\HVTLead{La reducción del consumo de agua requiere una red separada y una calidad compatible con el destino del agua tratada.}
\HVTText{La EPA estudia el reúso no potable a escala de edificio con tratamiento y evaluación de riesgos según fuente y uso. Esa evidencia orienta la formulación, pero no constituye una autorización de uso en Colombia. \HVTCite{3}}
\HVTText{El proyecto evaluará por separado las aguas de duchas y lavamanos; la inclusión de lavandería dependerá de su caracterización. Las aguas de inodoros tendrán una gestión independiente. La cocina no se incorporará automáticamente al circuito gris por su carga de grasas y residuos.}
\begin{HVTTable}{Etapa}{Diseño propuesto}{Punto de verificación}
\HVTTableRow{Separación}{Red identificada y sin conexión cruzada con agua potable.}{Planos y revisión de instalación.}
\HVTTableRow{Tratamiento}{Tecnología elegida según caudal, calidad y uso.}{Ensayos y desempeño bajo variación de carga.}
\HVTTableRow{Reúso}{Destino no potable definido y almacenamiento controlado.}{Calidad, señalización y operación mantenible.}
\end{HVTTable}
\HVTText{El riego de alimentos no se adoptará como destino predeterminado. El responsable sanitario y ambiental deberá establecer los usos admisibles, controles y requisitos aplicables al predio. Ante fallos de tratamiento se suspenderá el reúso y se activará la gestión alternativa prevista en el diseño.}
\end{HVTPage}
\begin{HVTPage}{Capítulo 8}{Energía}{Autonomía definida por el servicio a la familia}
\HVTLead{Primero se reduce la demanda; después se dimensionan la generación solar, el almacenamiento y el respaldo.}
\HVTText{Se levantará un perfil horario de cargas: iluminación, refrigeración, cocina, bombeo, tratamiento de agua, trabajo y movilidad. El balance distinguirá electricidad y calor. Los modelos fotovoltaicos y de baterías de SAM permiten estudiar producción y despacho temporal; sus resultados deben contrastarse con las condiciones del sitio. \HVTCite{4}}
\begin{HVTTable}{Alternativa}{Qué se evaluará}{Criterio de selección}
\HVTTableRow{Solar y red}{Autoconsumo, compras de energía y respaldo.}{Servicio, inversión y dependencia externa.}
\HVTTableRow{Solar y batería}{Potencia, capacidad útil, pérdidas y reposición.}{Horas de cobertura de cargas esenciales.}
\HVTTableRow{Hidrógeno}{Electrólisis, almacenamiento y reconversión separados.}{Balance completo, mantenimiento y viabilidad frente a alternativas.}
\end{HVTTable}
\HVTText{El almacenamiento mediante hidrógeno se plantea como una línea de evaluación progresiva, no como un equipo doméstico ya seleccionado. Se contará la energía de todos sus auxiliares y pérdidas. Su ubicación, protecciones y operación requerirán un diseño específico fuera de las áreas habitables.}
\HVTText{Un balance anual favorable no demuestra autonomía durante una noche o una semana de baja radiación. Se reportarán horas sin servicio, energía no suministrada y fracción de demanda cubierta, incluyendo escenarios de carga del vehículo sin asumir todavía su tecnología.}
\end{HVTPage}
\begin{HVTPage}{Capítulo 9}{Ciclos y materiales}{Producir alimentos y aprovechar los orgánicos}
\HVTLead{El paisaje productivo se dimensionará según el agua disponible, el suelo y el tiempo que la familia pueda dedicarle.}
\HVTText{La propuesta incluirá un huerto por etapas, especies adaptadas al lugar y cobertura del suelo. El registro separará producción cosechada, insumos, agua utilizada y trabajo de mantenimiento. El aporte alimentario se expresará mediante resultados del huerto, sin prometer autosuficiencia total.}
\HVTText{Los orgánicos se clasificarán en origen y se evaluará compostaje con manejo de humedad, olores y maduración. El biohidrógeno será una línea experimental separada, sujeta a caracterización de sustratos, balance de energía y validación de proceso; no se contará como suministro garantizado para la vivienda.}
\begin{HVTTable}{Componente}{Comparación propuesta}{Registro}
\HVTTableRow{Materiales}{Durabilidad, disponibilidad y mantenimiento.}{Origen, cantidades y reposiciones esperadas.}
\HVTTableRow{Construcción}{Movimiento de tierra y residuos.}{Volúmenes, destinos y medidas de reducción.}
\HVTTableRow{Huerto}{Producción y uso de recursos.}{Cosecha, agua e insumos por temporada.}
\HVTTableRow{Orgánicos}{Compostaje y alternativas de manejo.}{Masa de entrada y calidad del producto obtenido.}
\end{HVTTable}
\end{HVTPage}
\begin{HVTPage}{Capítulo 10}{Medir para decidir}{Un plan de evaluación con datos comparables}
\HVTLead{Cada promesa del proyecto tendrá una variable, un método de medida y un criterio de aceptación acordado antes de evaluar resultados.}
\begin{HVTTable}{Objetivo}{Indicador propuesto}{Comparación}
\HVTTableRow{Confort}{Horas ocupadas en condiciones aceptables.}{Alternativas con igual programa de uso.}
\HVTTableRow{Agua}{Volumen externo por habitante y calidad del reúso.}{Escenario de referencia y periodos de ocupación comparables.}
\HVTTableRow{Energía}{Demanda, cobertura y energía no suministrada.}{Perfiles horarios y condiciones climáticas documentadas.}
\HVTTableRow{Suelo}{Humedad, agua aplicada y condición del cultivo.}{Sectores y temporadas registrados.}
\end{HVTTable}
\HVTText{Los instrumentos compartirán fecha y hora; cada dato conservará unidad, ubicación y estado de calibración. Se definirán reglas para registros faltantes y valores atípicos antes del análisis. La información de hábitos familiares se limitará a lo necesario, con acceso acordado con los ocupantes.}
\HVTText{El estudio comenzará con simulaciones y una línea base. Después de construir, la puesta en marcha comprobará los sistemas por separado y finalmente en conjunto. Se contrastarán temporadas húmedas y secas; cualquier ahorro se comunicará junto con clima, ocupación y límites de la comparación.}
\end{HVTPage}
\begin{HVTPage}{Capítulo 11}{Desarrollo progresivo}{Etapas que terminan en decisiones concretas}
\begin{HVTTable}{Etapa}{Entregable}{Condición para avanzar}
\HVTTableRow{Definición}{Programa familiar, predio y presupuesto de estudio.}{Acuerdo de necesidades y responsables.}
\HVTTableRow{Diagnóstico}{Topografía, clima, suelo y restricciones.}{Aptitud del lugar y vacíos de información resueltos.}
\HVTTableRow{Anteproyecto}{Implantaciones, balances y alternativas de sistemas.}{Selección con costo y desempeño comparables.}
\HVTTableRow{Ingeniería}{Planos coordinados y especificaciones.}{Revisión profesional y permisos aplicables.}
\HVTTableRow{Ejecución}{Registros de obra y puesta en marcha.}{Pruebas funcionales y capacitación de la familia.}
\HVTTableRow{Seguimiento}{Informe de desempeño y ajustes.}{Resultados trazables durante periodos representativos.}
\end{HVTTable}
\HVTText{Se prevé integrar arquitectura, geotecnia, estructuras, instalaciones hidráulicas y eléctricas, ambiente y control. Los roles, honorarios y cronograma se acordarán cuando exista un predio definido. Las etapas permiten detener o simplificar una alternativa si su costo o mantenimiento no se justifican.}
\end{HVTPage}
\begin{HVTPage}{Capítulo 12}{Viabilidad}{Qué necesitamos para convertirlo en un piloto}
\HVTLead{El siguiente hito es seleccionar una familia y un predio para desarrollar el diagnóstico y el anteproyecto comparativo.}
\HVTText{El presupuesto distinguirá estudios, obra civil, sistemas de agua, generación, almacenamiento, instrumentación y mantenimiento. El análisis de ciclo de vida incorporará reposiciones y escenarios de precios; no se anunciará retorno económico sin cotizaciones y perfiles de uso verificables.}
\HVTText{Quedan por definir ubicación, área, habitantes, nivel de autonomía, servicios disponibles, presupuesto y aliados de ejecución. Esos datos cambiarán dimensiones y soluciones, pero no el método: conocer el lugar, reducir demanda, integrar sistemas y medir el desempeño. La casa se relaciona con los proyectos HVT de hogares eficientes y ecoaldea, conservando su escala unifamiliar y su propio registro experimental.}
\HVTHeading{Créditos y alcance de las imágenes}
\HVTText{Portada: Petra Nesti, Pexels 34533806. Patio: Amar Preciado, Pexels 15978605. Video: Juan Camilo Trujillo Botero, arquitectura colonial de Villa de Leyva, Pexels 37195962. Fotografías reducidas y video extractado sin audio; encuadre adaptado en la web. Licencia Pexels. No representan una obra ejecutada por HVT ni acreditan por sí mismos sostenibilidad. Las referencias sustentan fundamentos; las etapas y criterios descritos son una propuesta de trabajo de HVT. Edición: 10 de octubre de 2026.}
\end{HVTPage}
\begin{HVTPage}{Referencias}{Bibliografía}{Fuentes para desarrollar el diseño}
\begin{HVTReferences}
\HVTReference{Departamento de Energía de EE. UU. Consumer Guide to Passive Solar Home Design, 2021}{https://www.energy.gov/sites/default/files/2021-08/ES-Passive\%20Solar\%20design_081321.pdf}
\HVTReference{FAO. Irrigation Water Management: Introduction to Irrigation. Capítulo 2, Soil and Water}{https://www.fao.org/4/R4082E/r4082e03.htm}
\HVTReference{EPA. Onsite Non-Potable Water Reuse Research}{https://www.epa.gov/water-research/onsite-non-potable-water-reuse-research}
\HVTReference{System Advisor Model. PVWatts System Design: modelación fotovoltaica y almacenamiento}{https://samrepo.nlr.gov/help/pvwatts_system_design.html}
\HVTReference{Petra Nesti. Spanish Villa with Red Tiled Roofs and Courtyard. Pexels}{https://www.pexels.com/photo/spanish-villa-with-red-tiled-roofs-and-courtyard-34533806/}
\HVTReference{Pexels. Licencia gratuita para fotografías y videos}{https://www.pexels.com/license/}
\HVTReference{Amar Preciado. Courtyard of a Mexican Hacienda. Pexels}{https://www.pexels.com/photo/view-of-a-fountain-and-trees-in-the-courtyard-of-a-mexican-hacienda-15978605/}
\HVTReference{Juan Camilo Trujillo Botero. Aerial View of Historic Colombian Town. Pexels}{https://www.pexels.com/video/aerial-view-of-historic-colombian-town-37195962/}
\end{HVTReferences}
\HVTText{Consulta: octubre de 2026. Las guías internacionales son antecedentes técnicos, no sustituyen los estudios del predio ni los requisitos locales. El proyecto deberá documentar su matriz normativa con los profesionales responsables antes de la ingeniería de detalle.}
\end{HVTPage}
\end{document}
